% Detailed native teaching glyphs, reconstructed from the original Volume 3 art.
% Shared by this figure family; no title band, annotation box, shadow or texture.
\newcommand{\TBrain}[3]{\begin{scope}[shift={(#1,#2)},scale=#3]
 \path[draw=FigurePurple,fill=FigurePurpleFill,line width=.85pt,rounded corners=.4mm]
 (0,6.6) .. controls(-1.6,8.6)and(-4.2,8.4)..(-4.7,6.4) .. controls(-7.2,6.3)and(-8.1,4.6)..(-7.3,2.8) .. controls(-9.1,1.4)and(-8.9,-1.2)..(-7.4,-2.2) .. controls(-8.2,-4.4)and(-6.5,-6.7)..(-4.8,-6.4) .. controls(-4.1,-8.5)and(-1.4,-8.3)..(0,-6.8)
 .. controls(1.4,-8.3)and(4.1,-8.5)..(4.8,-6.4) .. controls(6.5,-6.7)and(8.2,-4.4)..(7.4,-2.2) .. controls(8.9,-1.2)and(9.1,1.4)..(7.3,2.8) .. controls(8.1,4.6)and(7.2,6.3)..(4.7,6.4) .. controls(4.2,8.4)and(1.6,8.6)..cycle;
 \draw[draw=FigurePurple,line width=.75pt] (0,6.6)--(0,-6.8);
 \draw[draw=FigurePurple,line width=1pt,line cap=round,line join=round] (-3.1,4.6)--(-3.1,1.1)--(-1.6,.1)--(-1.6,-5.5) (-5,1.3)--(-3.1,.2)--(-3.1,-2.5)--(-5,-3.7) (3.3,4.3)--(3.3,2.6)--(1.7,.7)--(1.7,-1) (5.4,1)--(3.7,-.5)--(3.7,-3.7)--(2.1,-5.1);
 \foreach \x/\y in {-3.1/4.6,-5/1.3,-5/-3.7,3.3/4.3,5.4/1,2.1/-5.1}{\draw[draw=FigurePurple,fill=white,line width=.8pt] (\x,\y) circle(.7);}
\end{scope}}
\newcommand{\TPerson}[5]{\begin{scope}[shift={(#1,#2)},scale=#3]
 \ifnum#5=1\def\pf{#4}\else\def\pf{white}\fi
 \draw[draw=#4,fill=\pf,line width=.65pt] (0,2) circle(1.55);
 \path[draw=#4,fill=\pf,line width=.65pt,rounded corners=.6pt] (-2.6,-3.5)--(-2.6,-1.4) arc[start angle=180,end angle=0,radius=2.6]--(2.6,-3.5)--cycle;
\end{scope}}
\newcommand{\TFolder}[5]{\begin{scope}[shift={(#1,#2)},scale=#3]
 \path[draw=FigureStroke,fill=#4Fill,line width=.75pt,rounded corners=.8mm] (-7,-4.5)--(-7,4.2)--(-3.5,4.2)--(-2,2.6)--(6,2.6)--(6,-4.5)--cycle;
 \ifnum#5=1\VEnvelope{0}{1.4}{.65}{FigureBlue}\fi
 \path[draw=FigureStroke,fill=#4Fill,line width=.75pt,rounded corners=.6mm] (-5.3,-4.5)--(-3.8,1)--(7,1)--(5.2,-4.5)--cycle;
\end{scope}}
\newcommand{\TBin}[5]{\begin{scope}[shift={(#1,#2)},scale=#3]
 \ifnum#5=1\begin{scope}[shift={(0,4)},rotate=-20]\VEnvelope{0}{0}{.7}{FigureBlue}\end{scope}\fi
 \path[draw=#4,fill=#4Fill,line width=.85pt,rounded corners=.7mm] (-4,3)--(-3.4,-6)--(3.4,-6)--(4,3)--cycle;
 \draw[draw=#4,line width=1.5pt,line cap=round] (-5,3)--(5,3);
 \draw[draw=#4,line width=.85pt,rounded corners=.5mm] (-1.8,3)--(-1.8,5)--(1.8,5)--(1.8,3);
 \foreach \x in {-2,0,2}{\draw[draw=#4,line width=.85pt,line cap=round] (\x,1)--(\x,-4.5);}
\end{scope}}
\newcommand{\TShield}[5]{\begin{scope}[shift={(#1,#2)},scale=#3]
 \path[draw=#4,fill=#4,line width=.65pt] (0,5.5) .. controls(1.9,4.2)and(3.6,3.8)..(4.4,3.5)--(4.4,-.4) .. controls(4.4,-3)and(2.2,-4.4)..(0,-5.7) .. controls(-2.2,-4.4)and(-4.4,-3)..(-4.4,-.4)--(-4.4,3.5) .. controls(-3.6,3.8)and(-1.9,4.2)..cycle;
 \ifnum#5=1\draw[draw=white,line width=1.2pt,line cap=round,line join=round] (-2.2,-.1)--(-.7,-1.8)--(2.6,1.7);\else\draw[draw=white,line width=1.2pt,line cap=round] (-1.7,-1.8)--(1.7,1.8) (-1.7,1.8)--(1.7,-1.8);\fi
\end{scope}}
\newcommand{\TTray}[4]{\begin{scope}[shift={(#1,#2)},scale=#3]
 \path[draw=#4,fill=#4Fill,line width=.85pt,rounded corners=.6mm] (-6,-3.3)--(-6,-.8)--(-3.8,3.3)--(3.8,3.3)--(6,-.8)--(6,-3.3)--cycle;
 \path[fill=#4,draw=#4,line width=.6pt,rounded corners=.4mm] (-6,-.3)--(-2.6,-.3)--(-1.9,-1.8)--(1.9,-1.8)--(2.6,-.3)--(6,-.3)--(6,-3.3)--(-6,-3.3)--cycle;
\end{scope}}
\newcommand{\TGear}[4]{\begin{scope}[shift={(#1,#2)},scale=#3]
 \foreach \a in {0,45,...,315}{\begin{scope}[rotate=\a]\path[draw=#4,fill=#4Fill,line width=.7pt,rounded corners=.4mm] (-1.5,3.3) rectangle (1.5,5.6);\end{scope}}
 \filldraw[draw=#4,fill=#4Fill,line width=.8pt] (0,0) circle(4.5);
 \filldraw[draw=#4,fill=white,line width=.8pt] (0,0) circle(2.1);
\end{scope}}
\newcommand{\TBalance}[3]{\begin{scope}[shift={(#1,#2)},scale=#3]
 \path[draw=FigureStroke,fill=FigureYellowFill,line width=.65pt,rounded corners=.35mm] (-1,-4.9)--(-.7,4.7)--(.7,4.7)--(1,-4.9)--(3,-5.5)--(3,-6)--(-3,-6)--(-3,-5.5)--cycle;
 \draw[draw=FigureYellow,line width=1.7pt,line cap=round] (-5.5,3.4)--(5.5,3.4);
 \filldraw[draw=FigureStroke,fill=FigureYellowFill,line width=.6pt] (0,4) circle(.8);
 \foreach \x in {-4,4}{\draw[draw=FigureStroke,line width=.6pt] (\x,3.4)--(\x-2.2,-1.2) (\x,3.4)--(\x+2.2,-1.2);\path[draw=FigureStroke,fill=FigureBlue,line width=.6pt] (\x-2.3,-1.2) arc[start angle=180,end angle=360,x radius=2.3,y radius=1.7]--cycle;}
\end{scope}}

\newcommand{\TOpenMail}[4]{\begin{scope}[shift={(#1,#2)},scale=#3,line join=round,line cap=round]
 \path[draw=FigureStroke,fill=#4Fill,line width=.75pt,rounded corners=.6mm] (-6,1)--(0,6.5)--(6,1)--(6,-5.2)--(-6,-5.2)--cycle;
 \draw[draw=FigureStroke,fill=white,line width=.7pt,rounded corners=.45mm] (-4.5,-1.5) rectangle (4.5,7.5);
 \foreach \y/\len in {5.5/5.5,3.5/5.5,1.5/4}{\draw[draw=#4,line width=.8pt] (-2.8,\y)--++(\len,0);}
 \path[draw=FigureStroke,fill=#4Fill,line width=.75pt,rounded corners=.6mm] (-6,1)--(0,-2.5)--(6,1)--(6,-5.2)--(-6,-5.2)--cycle;
 \draw[draw=FigureStroke,line width=.6pt] (-5.8,-5)--(-1.4,-1.7) (5.8,-5)--(1.4,-1.7);
\end{scope}}
% A uniform-tone variant preserves the reference neural glyph's layer count,
% filled circles and readable vertical spacing without replacing it by a logo.
\newcommand{\TNetwork}[7]{\begin{scope}[shift={(#1,#2)},scale=#3,yscale=1.15]
 \colorlet{VNetworkInput}{#7}\colorlet{VNetworkHidden}{#7}\colorlet{VNetworkOutput}{#7}
 \VLayerNet{0}{0}{1}{#4}{#5}{#6}
\end{scope}}
